\chapter{含参数积分}
\label{chap:parameter-integrals}

积分
\[
 F(t)=\int_X f(t,x)\,d\mu(x)
\]
把函数族压缩成一个参数函数。若要判断 \(F\) 是否可测、连续或可导，就要
把相应的极限移入积分号内。每一种结论都包含两个不同的问题：对固定参数，
内层积分是否有意义；当参数变化时，是否有一个可积函数控制这族被积函数。

\section{参数积分的可测性}

\begin{theorem}
\label{thm:parameter-integral-measurable}
设 \((T,\mathcal T,\nu)\)、\((X,\mathcal A,\mu)\) 为
\(\sigma\)-有限测度空间。
\begin{enumerate}[label=\textup{(\roman*)}]
 \item 若 \(f:T\times X\to[0,+\infty]\) 可测，则
 \[
  t\longmapsto\int_Xf(t,x)\,d\mu(x)
 \]
 是可测函数。
 \item 若 \(f\in L^1(\nu\times\mu)\)，则上述积分对几乎所有 \(t\)
 绝对收敛；把例外点上的值定义为零后得到 \(L^1(\nu)\) 函数，且
 \[
  \int_TF\,d\nu=\int_{T\times X}f\,d(\nu\times\mu).
 \]
\end{enumerate}
\end{theorem}

\begin{proof}
两项分别是 Tonelli 定理和 Fubini 定理。应注意：仅知道每个截面
\(x\mapsto f(t,x)\) 可测，并不自动推出 \(f\) 在乘积空间可测；联合可测
或其他足以保证它的条件必须另行验证。
\end{proof}

\begin{example}
若 \(E\subset T\times X\) 可测，则
\[
 F(t)=\mu(E_t)
\]
可测。这可用于描述几何区域的截面体积，也说明“截面长度随参数可测”不是
额外的光滑性结论，而是乘积可测性与 Tonelli 定理的结果。
\end{example}

\section{连续性}

\begin{theorem}[参数积分的连续性]
\label{thm:parameter-integral-continuity}
设 \((X,\mathcal A,\mu)\) 为测度空间，\(T\) 为度量空间，\(t_0\in T\)，
\(f:T\times X\to\mathbb R\)（或\(\mathbb C\)）的每个截面
\(f(t,\cdot)\) 都可测。假定：
\begin{enumerate}[label=\textup{(\roman*)}]
 \item 对几乎所有 \(x\)，\(f(t,x)\to f(t_0,x)\) 当 \(t\to t_0\)；
 \item 存在 \(t_0\) 的邻域 \(U\) 和 \(g\in L^1(\mu)\)，使
 \(|f(t,x)|\le g(x)\) 对 \(t\in U\) 几乎处处成立。
\end{enumerate}
则 \(F(t)=\int_Xf(t,x)\,d\mu(x)\) 在 \(U\) 上有定义，并在 \(t_0\) 连续。
\end{theorem}

\begin{proof}
任取 \(t_n\to t_0\)，舍去有限项后可令所有 \(t_n\in U\)。把这可数多个
参数对应的控制不等式的零测例外集合并，便可对整列使用同一个控制函数。
对 \(f(t_n,\cdot)\) 使用控制收敛定理，得到
\[
 F(t_n)\longrightarrow F(t_0).
\]
度量空间中序列连续等价于连续，故结论成立。
\end{proof}

控制条件只需在 \(t_0\) 的一个邻域成立。若 \(T\) 紧，并且每个点附近都
有一个控制函数，可取有限子覆盖处理整个参数集；一般不能把不同邻域上的
控制函数未经验证地拼成一个全局可积控制。

\begin{example}[Laplace 型积分]
对 \(t>0\)，令
\[
 F(t)=\int_0^\infty e^{-tx}\,dx.
\]
在固定 \(t_0>0\) 的邻域 \(t\ge t_0/2\) 内，
\(e^{-tx}\le e^{-t_0x/2}\in L^1(0,\infty)\)，故 \(F\) 在
\((0,\infty)\) 连续。这里没有一个对所有 \(t>0\) 都适用的可积控制
函数；局部控制已经足够。
\end{example}

\section{极限与积分}

若 \(t_n\to t_0\)，定理 \ref{thm:parameter-integral-continuity} 就是
控制收敛定理在参数族上的直接形式。还有两种常用替代条件：

\begin{proposition}
\label{prop:parameter-monotone-fatou}
设 \(f_t\ge0\) 可测。
\begin{enumerate}[label=\textup{(\roman*)}]
 \item 若 \(t_n\to t_0\) 且 \(f_{t_n}\uparrow f_{t_0}\) 几乎处处，则
 \(\int f_{t_n}\uparrow\int f_{t_0}\)。
 \item 对任意 \(t_n\to t_0\)，
 \[
  \int\liminf_n f_{t_n}\le\liminf_n\int f_{t_n}.
 \]
\end{enumerate}
\end{proposition}

第一项来自单调收敛，第二项来自 Fatou 引理。它们给出的分别是等式和单侧
估计。实际使用时，先把连续性问题转化为任意参数序列，再选择适合该函数族
的收敛定理。

\begin{counterexample}
令 \(f(t,x)=t^{-1}\mathbf1_{(0,t)}(x)\)（\(0<t\le1\)，\(0<x<1\)），并令
\(f(0,x)=0\)。对每个 \(x>0\)，当 \(t\downarrow0\) 时
\(f(t,x)\to0\)，但
\[
 \int_0^1f(t,x)\,dx=1.
\]
尖峰高度与支撑缩小恰好抵消；不存在邻域内的共同可积控制函数。
\end{counterexample}

\section{积分号下求导}

\begin{theorem}[积分号下求导]
\label{thm:differentiation-under-integral}
设 \((X,\mathcal A,\mu)\) 为测度空间，\(I\subset\mathbb R\) 为开区间，
\(f:I\times X\to\mathbb R\)，\(t_0\in I\)。假定存在含 \(t_0\) 的
开区间 \(U\subset I\)，使：
\begin{enumerate}[label=\textup{(\roman*)}]
 \item 对每个 \(t\in U\)，\(f(t,\cdot)\) 可测；
 \item 存在一个零测集 \(N\)，使对每个 \(x\notin N\)，函数
 \(t\mapsto f(t,x)\) 在整个 \(U\) 上可导；
 \item 存在 \(g\in L^1(\mu)\)，使对每个 \(x\notin N\) 和每个 \(t\in U\)，
 \[
  |\partial_tf(t,x)|\le g(x)
 \]；
 \item \(f(t_0,\cdot)\in L^1(\mu)\)。
\end{enumerate}
在 \(N\) 上把偏导数定义为零。则 \(F(t)=\int_Xf(t,x)\,d\mu(x)\)
在 \(U\) 上有定义，在 \(t_0\) 可导，且
\[
 F'(t_0)=\int_X\partial_tf(t_0,x)\,d\mu(x).
\]
\end{theorem}

\begin{proof}
先由中值定理得到
\(|f(t,x)|\le |f(t_0,x)|+|t-t_0|g(x)\)（\(x\notin N\)），
故每个 \(t\in U\) 的截面都可积。再取任意非零实数列 \(h_n\to0\)，
使连接 \(t_0\) 与 \(t_0+h_n\) 的整段都包含在 \(U\) 内。
这些可测差商在 \(X\setminus N\) 上的极限是偏导数，因而补零后的偏导数
可测。对足够小的非零 \(h\)，固定
\(x\notin N\)，对函数 \(t\mapsto f(t,x)\) 使用一元中值定理，得到
\[
 \left|\frac{f(t_0+h,x)-f(t_0,x)}{h}\right|\le g(x).
\]
这里同一个零测集 \(N\) 对所有小 \(h\) 都适用；这正是条件 (ii)--(iii)
中量词写法的作用。差商在 \(X\setminus N\) 上逐点趋于
\(\partial_tf(t_0,x)\)。对任意这样的序列 \(h_n\) 使用控制收敛，
再用实变量极限的序列判据，得到
\[
 \lim_{h\to0}\frac{F(t_0+h)-F(t_0)}h
 =\int_X\lim_{h\to0}
 \frac{f(t_0+h,x)-f(t_0,x)}h\,d\mu(x),
\]
即所求公式。复值函数可分别对实部和虚部应用这一结论；
不需要对复值函数使用实中值定理。
\end{proof}

\begin{remark}
条件 (iii) 可以改成直接控制差商，而不要求偏导数在整个邻域都存在：
若
\[
 \left|\frac{f(t_0+h,x)-f(t_0,x)}h\right|\le g(x)
\]
对所有足够小的 \(h\) 成立，且差商几乎处处收敛，同一证明仍适用。这是
定理中实际承担极限交换工作的条件。
\end{remark}

\begin{corollary}[Leibniz 公式]
\label{cor:leibniz-variable-limits}
设 \(\alpha,\beta\) 在 \(t_0\) 可导。假定存在 \(t_0\) 的邻域 \(U\) 和
有界区间 \(J\)，使 \(\alpha(t),\beta(t)\in J\) 对 \(t\in U\) 成立，
并且 \(f:U\times J\to\mathbb R\) 满足：
\begin{enumerate}[label=\textup{(\roman*)}]
 \item \(f(t,\cdot)\) 可测，且 \(f(t_0,\cdot)\in L^1(J)\)；
 \item 存在零测集 \(N\subset J\)，使对每个 \(x\in J\setminus N\)，
 \(\partial_tf(t,x)\) 在全部 \(t\in U\) 上存在，并有 \(g\in L^1(J)\)
 满足 \(|\partial_tf(t,x)|\le g(x)\) 对这些 \((t,x)\) 同时成立；
 \item \(f\) 在
 \((t_0,\alpha(t_0))\) 与 \((t_0,\beta(t_0))\) 处联合连续。
\end{enumerate}
则
\[
 \frac d{dt}\int_{\alpha(t)}^{\beta(t)}f(t,x)\,dx\bigg|_{t=t_0}
 =
 f(t_0,\beta(t_0))\beta'(t_0)
 -f(t_0,\alpha(t_0))\alpha'(t_0)
 +\int_{\alpha(t_0)}^{\beta(t_0)}
   \partial_tf(t_0,x)\,dx.
\]
\end{corollary}

\begin{proof}
固定 \(c\in J\)，先考察
\[
 H_\beta(t)=\int_c^{\beta(t)}f(t,x)\,dx.
\]
记 \(\beta_0=\beta(t_0)\)、\(\beta_h=\beta(t_0+h)\)。差可分成
\[
 \begin{split}
 H_\beta(t_0+h)-H_\beta(t_0)
 ={}&\int_c^{\beta_0}
       \bigl(f(t_0+h,x)-f(t_0,x)\bigr)\,dx\\
 &+\int_{\beta_0}^{\beta_h}f(t_0+h,x)\,dx.
 \end{split}
\]
第一项除以 \(h\) 后，由积分号下求导定理趋于
\(\int_c^{\beta_0}\partial_tf(t_0,x)\,dx\)。若 \(\beta_h\ne\beta_0\)，
第二项的差商写成
\[
 \frac{\beta_h-\beta_0}{h}\,
 \frac1{\beta_h-\beta_0}
 \int_{\beta_0}^{\beta_h}f(t_0+h,x)\,dx.
\]
联合连续性使第二个因子趋于 \(f(t_0,\beta_0)\)，第一个因子趋于
\(\beta'(t_0)\)；若端点在某子列上不动，结论仍成立。因此
\[
 H_\beta'(t_0)
 =\int_c^{\beta_0}\partial_tf(t_0,x)\,dx
  +f(t_0,\beta_0)\beta'(t_0).
\]
对 \(H_\alpha(t)=\int_c^{\alpha(t)}f(t,x)\,dx\) 作同一计算，再用原积分
等于 \(H_\beta-H_\alpha\)，便得到公式。
\end{proof}

\begin{example}
对 \(t>0\)，
\[
 \Gamma(t)=\int_0^\infty x^{t-1}e^{-x}\,dx.
\]
若 \(t\) 限制在紧子区间 \([\alpha,\beta]\subset(0,\infty)\)，则
\[
 |x^{t-1}(\log x)e^{-x}|
\le
\begin{cases}
 x^{\alpha-1}|\log x|,&0<x\le1,\\
 x^{\beta-1}(\log x)e^{-x},&x>1,
\end{cases}
\]
右侧可积。因此
\[
 \Gamma'(t)=\int_0^\infty x^{t-1}(\log x)e^{-x}\,dx.
\]
对任意整数 \(k\ge0\)，把上述控制函数中的 \(|\log x|\) 改为
\(|\log x|^k\)。在零附近作代换 \(x=e^{-u}\)，有
\[
 \int_0^1x^{\alpha-1}|\log x|^k\,dx
 =\int_0^\infty e^{-\alpha u}u^k\,du<\infty.
\]
无穷远处，\(x^{\beta-1}(\log x)^ke^{-x}\le C e^{-x/2}\)
对充分大的 \(x\) 成立。因而每一阶导数都有局部可积控制。
依次应用求导定理得到
\(\Gamma^{(k)}(t)=\int_0^\infty x^{t-1}(\log x)^ke^{-x}\,dx\)，
再用连续性定理说明每个导数连续，故 \(\Gamma\in C^\infty(0,\infty)\)。
\end{example}

\section*{习题}
\addcontentsline{toc}{section}{习题}
\subsection*{A组　可测性与连续性}
\begin{enumerate}[label=\textbf{\thechapter.\arabic*.}]
\exerciseitem{parameter-section-measurability}
设 \(T,X\) 为本章可测性定理中的两个\(\sigma\)-有限测度空间，
\(f:T\times X\to\mathbb R\) 联合可测。证明
\[
 t\mapsto\mu(\{x:f(t,x)>a\})
\]
对每个 \(a\in\mathbb R\) 可测。

\exerciseitem{parameter-continuity-sequential}
补全定理 \ref{thm:parameter-integral-continuity} 中从序列连续到连续的
论证，并说明若参数空间不是第一可数空间，为什么该证明需要调整。

\exerciseitem{laplace-continuity-differentiability}
设 \(h:(0,\infty)\to\mathbb R\) 可测，\(a\in\mathbb R\)，且
\(e^{-ax}|h(x)|\) 可积。证明
\[
 F(t)=\int_0^\infty e^{-tx}h(x)\,dx
\]
在 \(t>a\) 连续；若还假设 \(xe^{-ax}|h(x)|\) 可积，
则证明它在 \(t>a\) 可导并求导。

\end{enumerate}
\subsection*{B组　积分号下求导}
\begin{enumerate}[label=\textbf{\thechapter.\arabic*.},resume]
\exerciseitem{parameter-no-domination}
完整验证正文尖峰反例，并证明任何控制所有
\(0<t<t_0\) 的函数 \(g\) 都不可能属于 \(L^1(0,1)\)。

\exerciseitem{differentiate-exponential}
对 \(t>0\)，在积分号下求导计算
\[
 \int_0^\infty e^{-tx}\,dx,\qquad
 \int_0^\infty xe^{-tx}\,dx.
\]

\exerciseitem{differentiate-gaussian}
设 \(F(t)=\int_{\mathbb R}e^{-x^2}\cos(tx)\,dx\)。证明
\(F'(t)=-\int xe^{-x^2}\sin(tx)\,dx\)，并继续求 \(F''\)。

\exerciseitem{leibniz-rule-proof}
从固定积分区间的求导定理和微积分基本定理完整推出推论
\ref{cor:leibniz-variable-limits}。

\end{enumerate}
\subsection*{C组　推广与应用}
\begin{enumerate}[label=\textbf{\thechapter.\arabic*.},resume]
\exerciseitem{log-integral-parameter}
对 \(t>-1\)，令
\[
 F(t)=\int_0^1\frac{x^t-1}{\log x}\,dx,
\]
在适当解释 \(x=1\) 处取值后证明 \(F'(t)=1/(t+1)\)，并求 \(F(t)\)。

\optionalexerciseitem{holomorphic-parameter-integral}
本题在学过 Cauchy 积分公式后完成。设 \((X,\mathcal A,\mu)\) 为测度空间，
\(D\subset\mathbb C\) 开，\(f:D\times X\to\mathbb C\) 的每个
\(x\)-截面可测；存在公共零测集 \(N\)，对 \(x\notin N\)，
\(f(\cdot,x)\) 全纯。对每个紧集 \(K\Subset D\)，有
\(g_K\in L^1(\mu)\) 使 \(|f(z,x)|\le g_K(x)\) 对所有 \(z\in K\)
及公共零测例外集外的 \(x\) 成立。利用 Cauchy 导数估计及控制收敛
证明 \(F(z)=\int f(z,x)\,d\mu(x)\) 全纯。

\projectexerciseitem{parameter-integral-checklist}
选择一个概率论中的矩母函数或物理中的配分函数，逐项核验可测性、积分有限、
连续性和积分号下求导条件；不能只形式地写出求导结果。
\end{enumerate}
